\section{Energy--momentum tensor constructed from the flowed
fields}
\label{sec:4}
\subsection{Small flow-time expansion and the renormalization group
argument}
To express the energy--momentum tensor in Eqs.~\eqref{eq:(2.7)}
and~\eqref{eq:(2.10)} in terms of the flowed fields, we consider following
second-rank symmetric tensors (which are even under the CP transformation)
constructed from the flowed fields:
\begin{align}
   \Tilde{\mathcal{O}}_{1\mu\nu}(t,x)&\equiv
   G_{\mu\rho}^a(t,x)G_{\nu\rho}^a(t,x),
\label{eq:(4.1)}\\
   \Tilde{\mathcal{O}}_{2\mu\nu}(t,x)&\equiv
   \delta_{\mu\nu}G_{\rho\sigma}^a(t,x)G_{\rho\sigma}^a(t,x),
\label{eq:(4.2)}\\
   \Tilde{\mathcal{O}}_{3\mu\nu}(t,x)&\equiv
   \mathring{\Bar{\chi}}(t,x)
   \left(\gamma_\mu\overleftrightarrow{D}_\nu
   +\gamma_\nu\overleftrightarrow{D}_\mu\right)
   \mathring{\chi}(t,x),
\label{eq:(4.3)}\\
   \Tilde{\mathcal{O}}_{4\mu\nu}(t,x)&\equiv
   \delta_{\mu\nu}
   \mathring{\Bar{\chi}}(t,x)
   \overleftrightarrow{\Slash{D}}
   \mathring{\chi}(t,x),
\label{eq:(4.4)}\\
   \Tilde{\mathcal{O}}_{5\mu\nu}(t,x)&\equiv
   \delta_{\mu\nu}
   m\mathring{\Bar{\chi}}(t,x)
   \mathring{\chi}(t,x).
\label{eq:(4.5)}
\end{align}
Note that all the above operators $\Tilde{\mathcal{O}}_{i\mu\nu}(t,x)$ are of
dimension~$4$ for any~$D$.

We also introduce corresponding bare operators in the $D$-dimensional
$x$-space:
\begin{align}
   \mathcal{O}_{1\mu\nu}(x)&\equiv
   F_{\mu\rho}^a(x)F_{\nu\rho}^a(x),
\label{eq:(4.6)}\\
   \mathcal{O}_{2\mu\nu}(x)&\equiv
   \delta_{\mu\nu}F_{\rho\sigma}^a(x)F_{\rho\sigma}^a(x),
\label{eq:(4.7)}\\
   \mathcal{O}_{3\mu\nu}(x)&\equiv
   \Bar{\psi}(x)
   \left(\gamma_\mu\overleftrightarrow{D}_\nu
   +\gamma_\nu\overleftrightarrow{D}_\mu\right)
   \psi(x),
\label{eq:(4.8)}\\
   \mathcal{O}_{4\mu\nu}(x)&\equiv
   \delta_{\mu\nu}
   \Bar{\psi}(x)
   \overleftrightarrow{\Slash{D}}
   \psi(x),
\label{eq:(4.9)}\\
   \mathcal{O}_{5\mu\nu}(x)&\equiv
   \delta_{\mu\nu}
   m_0\Bar{\psi}(x)
   \psi(x).
\label{eq:(4.10)}
\end{align}
The mass dimension of $\mathcal{O}_{1\mu\nu}(x)$ and~$\mathcal{O}_{2\mu\nu}(x)$
is~$4$, while that of $\mathcal{O}_{3\mu\nu}(x)$, $\mathcal{O}_{4\mu\nu}(x)$,
and~$\mathcal{O}_{5\mu\nu}(x)$ is~$D$.

We now consider the situation in which the flow time~$t$
in~Eqs.~\eqref{eq:(4.1)}--\eqref{eq:(4.5)} is very small. Since a flowed field
at position~$x$ with a flow time~$t$ is a combination of un-flowed fields in
the vicinity of~$x$ of radius~$\sim\sqrt{8t}$,\footnote{This follows from
the fact that the flow equations are basically the diffusion equation.} the
operators~\eqref{eq:(4.1)}--\eqref{eq:(4.5)} can be regarded as local operators
in $x$-space in the $t\to0$ limit. Since the bare operators
in~Eqs.~\eqref{eq:(4.6)}--\eqref{eq:(4.10)} span a complete set of symmetric
second-rank gauge-invariant local operators of dimension~$4$ (for $D\to4$)
which are even under CP, we have the following (asymptotic) expansion for
small~$t$:
\begin{equation}
   \Tilde{\mathcal{O}}_{i\mu\nu}(t,x)
   =\left\langle\Tilde{\mathcal{O}}_{i\mu\nu}(t,x)\right\rangle
   +\zeta_{ij}(t)
   \left[\mathcal{O}_{j\mu\nu}(x)
   -\left\langle\mathcal{O}_{j\mu\nu}(x)\right\rangle\right]+O(t),
\label{eq:(4.11)}
\end{equation}
where the abbreviated terms are contributions of operators of mass dimension
$6$ (for $D\to4$) or higher.

In writing down the small flow-time expansion~\eqref{eq:(4.11)}, we have
assumed that there is no other $D$-dimensional composite operator at the
point~$x$. That the expansion~\eqref{eq:(4.11)} cannot necessarily hold when
there is another operator at the point~$x$ [say, $\mathcal{P}(x)$] can be seen
by noting that the product of the left-hand side of Eq.~\eqref{eq:(4.11)}
with~$\mathcal{P}(x)$ does not possess any divergence for~$t>0$, while each
term in the right-hand side can have an equal-point singularity
with~$\mathcal{P}(x)$ because the operators in the right-hand side
of~Eq.~\eqref{eq:(4.11)} are $D$-dimensional (i.e., non-flowed) composite
operators. This is a contradiction if Eq.~\eqref{eq:(4.11)} holds. This implies
that the formula we will derive for the energy--momentum tensor below holds
\emph{only when the energy--momentum tensor has no overlap with other
operators}.\footnote{This important point was not fully recognized
in~Ref.~\cite{Suzuki:2013gza}.} In particular, we cannot say anything about
whether the Ward--Takahashi relation~\eqref{eq:(2.9)} is reproduced with our
construction. Still, our construction is expected to have a correct
normalization because it is determined from matching the energy--momentum
tensor in the dimensional regularization which fulfills~Eq.~\eqref{eq:(2.9)}.
Our lattice energy--momentum tensor is thus useful to compute correlation
functions in which the energy--momentum tensor is separated from other
operators. This is the case for correlation functions relevant to the
viscosities~\cite{Nakamura:2004sy,Meyer:2007ic,Meyer:2007dy}, for example.

The expansion~\eqref{eq:(4.11)} may be inverted as
\begin{equation}
   \mathcal{O}_{i\mu\nu}(x)-\left\langle\mathcal{O}_{i\mu\nu}(x)\right\rangle
   =\left(\zeta^{-1}\right)_{ij}(t)
   \left[\Tilde{\mathcal{O}}_{j\mu\nu}(t,x)
   -\left\langle\Tilde{\mathcal{O}}_{j\mu\nu}(t,x)\right\rangle\right]
   +O(t),
\label{eq:(4.12)}
\end{equation}
where $\zeta^{-1}$ denotes the inverse matrix of~$\zeta$. Then, by substituting
this relation into the energy--momentum tensor~\eqref{eq:(2.7)}
and~\eqref{eq:(2.10)} in terms of the bare operators,
\begin{align}
   \left\{T_{\mu\nu}\right\}_R(x)
   &=\frac{1}{g_0^2}\left\{
   \mathcal{O}_{1\mu\nu}(x)
   -\left\langle\mathcal{O}_{1\mu\nu}(x)\right\rangle
   -\frac{1}{4}
   \left[\mathcal{O}_{2\mu\nu}(x)
   -\left\langle\mathcal{O}_{2\mu\nu}(x)\right\rangle\right]
   \right\}
\notag\\
   &\qquad{}
   +\frac{1}{4}
   \left[\mathcal{O}_{3\mu\nu}(x)
   -\left\langle\mathcal{O}_{3\mu\nu}(x)\right\rangle\right]
   -\frac{1}{2}
   \left[\mathcal{O}_{4\mu\nu}(x)
   -\left\langle\mathcal{O}_{4\mu\nu}(x)\right\rangle\right]
\notag\\
   &\qquad\qquad{}
   -\left[\mathcal{O}_{5\mu\nu}(x)
   -\left\langle\mathcal{O}_{5\mu\nu}(x)\right\rangle\right],
\label{eq:(4.13)}
\end{align}
we have the expression (for $D=4$)
\begin{align}
   \left\{T_{\mu\nu}\right\}_R(x)
   &=c_1(t)\left[
   \Tilde{\mathcal{O}}_{1\mu\nu}(t,x)
   -\frac{1}{4}\Tilde{\mathcal{O}}_{2\mu\nu}(t,x)
   \right]
\notag\\
   &\qquad{}
   +c_2(t)\left[
   \Tilde{\mathcal{O}}_{2\mu\nu}(t,x)
   -\left\langle\Tilde{\mathcal{O}}_{2\mu\nu}(t,x)\right\rangle
   \right]
\notag\\
   &\qquad\qquad{}
   +c_3(t)\left[
   \Tilde{\mathcal{O}}_{3\mu\nu}(t,x)
   -2\Tilde{\mathcal{O}}_{4\mu\nu}(t,x)
   -\left\langle
   \Tilde{\mathcal{O}}_{3\mu\nu}(t,x)
   -2\Tilde{\mathcal{O}}_{4\mu\nu}(t,x)
   \right\rangle
   \right]
\notag\\
   &\qquad\qquad\qquad{}
   +c_4(t)\left[
   \Tilde{\mathcal{O}}_{4\mu\nu}(t,x)
   -\left\langle\Tilde{\mathcal{O}}_{4\mu\nu}(t,x)\right\rangle
   \right]
\notag\\
   &\qquad\qquad\qquad\qquad{}
   +c_5(t)\left[
   \Tilde{\mathcal{O}}_{5\mu\nu}(t,x)
   -\left\langle\Tilde{\mathcal{O}}_{5\mu\nu}(t,x)\right\rangle
   \right]+O(t),
\label{eq:(4.14)}
\end{align}
where
\begin{align}
   &c_1(t)=\Tilde{c}_1(t),\qquad
   c_2(t)=\Tilde{c}_2(t)+\frac{1}{4}c_1(t),
\notag\\
   &c_3(t)=\Tilde{c}_3(t),\qquad
   c_4(t)=\Tilde{c}_4(t)+2c_3(t),\qquad
   c_5(t)=\Tilde{c}_5(t),
\label{eq:(4.15)}
\end{align}
and
\begin{equation}
   \Tilde{c}_i(t)\equiv\frac{1}{g_0^2}\left\{
   \left(\zeta^{-1}\right)_{1i}(t)
   -\frac{1}{4}\left(\zeta^{-1}\right)_{2i}(t)\right\}
   +\frac{1}{4}\left(\zeta^{-1}\right)_{3i}(t)
   -\frac{1}{2}\left(\zeta^{-1}\right)_{4i}(t)
   -\left(\zeta^{-1}\right)_{5i}(t).
\label{eq:(4.16)}
\end{equation}
In Eq.~\eqref{eq:(4.14)}, we have used the fact that the finite operator
$\Tilde{\mathcal{O}}_{1\mu\nu}(t,x)-(1/4)\Tilde{\mathcal{O}}_{2\mu\nu}(t,x)$ is
traceless in~$D=4$ and thus has no vacuum expectation value.
Equation~\eqref{eq:(4.14)} shows that if one knows the $t\to0$ behavior of the
coefficients~$c_i(t)$, the energy--momentum tensor can be obtained as the
$t\to0$ limit of the combination in the right-hand side. As already noted,
since the composite operators~\eqref{eq:(4.1)}--\eqref{eq:(4.5)} constructed
from (ringed) flowed fields should be independent of the regularization
adopted, one may use the lattice regularization to compute correlation
functions of the quantity in the right-hand side of~Eq.~\eqref{eq:(4.14)}. This
provides a possible method to compute correlation functions of the
correctly normalized conserved energy--momentum tensor with the lattice
regularization.

Thus, we are interested in the $t\to0$ behavior of the coefficients~$c_i(t)$
in~Eq.~\eqref{eq:(4.14)}. Quite interestingly, one can argue that the
coefficients~$c_i(t)$ can be evaluated by the perturbation theory for~$t\to0$
thanks to the asymptotic freedom. To see this, we apply
\begin{equation}
   \left(\mu\frac{\partial}{\partial\mu}\right)_0
\label{eq:(4.17)}
\end{equation}
on both sides of~Eq.~\eqref{eq:(4.14)}, where $\mu$ is the renormalization
scale and the subscript~$0$ implies that the derivative is taken while all bare
quantities are kept fixed. Since the energy--momentum tensor is not
multiplicatively renormalized,
$(\mu\partial/\partial\mu)_0(\text{left-hand side of Eq.~\eqref{eq:(4.14)}})=0$.
On the right-hand side, since $\Tilde{O}_{1,2,3,4\mu\nu}(t,x)$
and~$(1/m)\Tilde{O}_{5\mu\nu}(t,x)$ are entirely given by bare quantities
through the flow equations, we have
\begin{equation}
   \left(\mu\frac{\partial}{\partial\mu}\right)_0\Tilde{O}_{1,2,3,4\mu\nu}(t,x)
   =\left(\mu\frac{\partial}{\partial\mu}\right)_0
   \frac{1}{m}\Tilde{O}_{5\mu\nu}(t,x)
   =0.
\label{eq:(4.18)}
\end{equation}
These observations imply
\begin{equation}
   \left(\mu\frac{\partial}{\partial\mu}\right)_0c_{1,2,3,4}(t)
   =\left(\mu\frac{\partial}{\partial\mu}\right)_0mc_5(t)
   =0.
\label{eq:(4.19)}
\end{equation}
Then the standard renormalization group argument says that $c_{1,2,3,4}(t)$
and~$mc_5(t)$ are independent of the renormalization scale, if the renormalized
parameters in these quantities are replaced by running parameters defined by
\begin{align}
   &q\frac{\mathrm{d}\Bar{g}(q)}{\mathrm{d}q}=\beta(\Bar{g}(q)),\qquad
   \Bar{g}(q=\mu)=g,
\label{eq:(4.20)}\\
   &q\frac{\mathrm{d}\Bar{m}(q)}{\mathrm{d}q}=-\gamma_m(\Bar{g}(q))\Bar{m}(q),\qquad
   \Bar{m}(q=\mu)=m,
\label{eq:(4.21)}
\end{align}
where $\mu$ is the original renormalization scale. Thus, since $c_{1,2,3,4}(t)$
and~$mc_5(t)$ are independent of the renormalization scale, two possible
choices, $q=\mu$ and~$q=1/\sqrt{8t}$, should give an identical result. In this
way, we infer that
\begin{align}
   c_{1,2,3,4}(t)(g,m;\mu)
   &=c_{1,2,3,4}(t)(\Bar{g}(1/\sqrt{8t}),\Bar{m}(1/\sqrt{8t});1/\sqrt{8t}),
\label{eq:(4.22)}\\
   c_5(t)(g,m;\mu)
   &=\frac{\Bar{m}(1/\sqrt{8t})}{m}
   c_5(t)(\Bar{g}(1/\sqrt{8t}),\Bar{m}(1/\sqrt{8t});1/\sqrt{8t}),
\label{eq:(4.23)}
\end{align}
where we have explicitly written the dependence of $c_i(t)$ on renormalized
parameters and on the renormalization scale. Finally, since the running gauge
coupling $\Bar{g}(1/\sqrt{8t})\to0$ for~$t\to0$ thanks to the asymptotic
freedom, we expect that we can compute $c_i(t)$ for~$t\to0$ by using the
perturbation theory. Although we are interested in low-energy physics for which
the perturbation theory is ineffective, the coefficients~$c_i(t)$
in~Eq.~\eqref{eq:(4.14)} for~$t\to0$ can be evaluated by perturbation theory;
this might be regarded as a sort of factorization.

\subsection{$c_i(t)$ to the one-loop order}
We thus evaluate the above coefficients~$c_i(t)$ in~Eq.~\eqref{eq:(4.14)} to
the one-loop order approximation. For this, we compute the mixing coefficients
$\zeta_{ij}(t)$ in~Eq.~\eqref{eq:(4.11)} to the one-loop order. Then $c_i(t)$
are obtained by~Eqs.~\eqref{eq:(4.15)} and~\eqref{eq:(4.16)}. The loop
expansion of~$\zeta_{ij}(t)$ would yield
\begin{equation}
   \zeta_{ij}(t)=\delta_{ij}+\zeta_{ij}^{(1)}(t)+\zeta_{ij}^{(2)}(t)+\dotsb,
\label{eq:(4.24)}
\end{equation}
for $j=1$ and~$j=2$, where the superscript denotes the loop order, and for
$j=3$, $4$, and~$5$, by taking the factor in~Eq.~\eqref{eq:(3.26)} into
account, we set 
% for $i=1$ and~$i=2$, where the superscript denotes the loop order, and for
% $i=3$, $4$, and~$5$, by taking the factor in~Eq.~\eqref{eq:(3.26)} into
% account, we set 
\begin{equation}
   \zeta_{ij}(t)=Z(\epsilon)
   \left[\delta_{ij}+\zeta_{ij}^{(1)}(t)+\zeta_{ij}^{(2)}(t)+\dotsb\right].
\label{eq:(4.25)}
\end{equation}

As Ref.~\cite{Suzuki:2013gza}, it is straightforward to compute
$\zeta_{ij}^{(1)}(t)$.\footnote{The justification for the following
computational prescription was not well explained in previous versions of the
present paper: One may wonder why the one-loop matching
coefficients~$\zeta_{ij}^{(1)}(t)$ can be read off from the correlation
functions~\eqref{eq:(4.26)} and~\eqref{eq:(4.29)} alone, without computing
corresponding correlation functions in which flowed composite
operators~$\Tilde{\mathcal{O}}_{i\mu\nu}(t,x)$ are replaced by bare
ones~$\mathcal{O}_{i\mu\nu}(x)$. The justification is directly related to our
way of treatment of infrared (IR) divergences and we supplement detailed
explanation in~Appendix~\ref{sec:D}. We are quite grateful to a referee of the
present paper for suggestions on this point.} For example, to obtain
$\zeta_{ij}^{(1)}(t)$ with~$j=1$ and~$2$, we consider the correlation function
\begin{equation}
   \left\langle\Tilde{\mathcal{O}}_{i\mu\nu}(t,x)
   A_\beta^b(y)A_\gamma^c(z)\right\rangle.
\label{eq:(4.26)}
\end{equation}
In the Feynman gauge which we use throughout the present paper, this has the
structure
\begin{equation}
   g_0^2\int_k\frac{\mathrm{e}^{ik(x-y)}}{k^2}g_0^2\int_\ell\frac{\mathrm{e}^{i\ell(x-z)}}{\ell^2}
   \delta^{bc}\mathcal{M}_{\mu\nu,\beta\gamma}(k,\ell).
\label{eq:(4.27)}
\end{equation}
After expanding $\mathcal{M}_{\mu\nu,\beta\gamma}(k,\ell)$ to~$O(k,\ell)$, we can
make use of the following correspondence to read off the operator mixing:
\begin{equation}
   ik_\mu i\ell_\nu\delta_{\beta\gamma}\to F_{\mu\rho}^a(x)F_{\nu\rho}^a(x),\qquad
   ik\cdot i\ell\delta_{\mu\nu}\delta_{\beta\gamma}
   \to\frac{1}{4}\delta_{\mu\nu}F_{\rho\sigma}^a(x)F_{\rho\sigma}^a(x).
\label{eq:(4.28)}
\end{equation}
In this way, we obtain~$\zeta_{ij}^{(1)}(t)$ with $j=1$ and~$2$.\footnote{For
the momentum integration, we use the integration formulas
in~Appendix~\ref{sec:B}.}

Similarly, to obtain $\zeta_{ij}^{(1)}(t)$ with $j=3$, $4$, and~$5$, we consider
\begin{equation}
   \left\langle\Tilde{\mathcal{O}}_{i\mu\nu}(t,x)
   \psi(y)\Bar{\psi}(z)\right\rangle,
\label{eq:(4.29)}
\end{equation}
whose general structure reads
\begin{equation}
   \int_k\frac{\mathrm{e}^{ik(y-x)}}{i\Slash{k}+m_0}
   \mathcal{M}_{\mu\nu}(k,\ell)
   \int_\ell\frac{\mathrm{e}^{i\ell(x-z)}}{i\Slash{\ell}+m_0}.
\label{eq:(4.30)}
\end{equation}
We then expand $\mathcal{M}_{\mu\nu}(k,\ell)$ to~$O(k)$ and~$O(\ell)$ and use the
correspondence
\begin{equation}
   \gamma_\mu i(k+\ell)_\nu+\gamma_\nu i(k+\ell)_\mu\to
   \Bar{\psi}(x)\left[
   \gamma_\mu\overleftrightarrow{D}_\nu+\gamma_\nu\overleftrightarrow{D}_\mu
   \right]\psi(x),\qquad
   \delta_{\mu\nu}\to\delta_{\mu\nu}\Bar{\psi}(x)\psi(x),
\label{eq:(4.31)}
\end{equation}
to read off the operator mixing.

% A remark on the correspondences in~Eqs.~\eqref{eq:(4.28)}
% and~\eqref{eq:(4.31)}: In our small flow-time expansion~\eqref{eq:(4.11)}, the
% operators~$\mathcal{O}_{j\mu\nu}(x)$ in the right-hand side are bare, not
% renormalized. Such bare operators in the momentum space are represented by
% tree-level vertices in the left-hand sides of~Eqs.~\eqref{eq:(4.28)}
% and~\eqref{eq:(4.31)}. This is in accord with the definition of the bare
% operator in the conventional operator renormalization, in which multiplicative
% renormalization factors to bare operators are chosen so that the insertion of
% corresponding tree-level vertices with the renormalization factors produces
% finite correlation functions.

% In the above calculation, we expand the integrand of the loop integration with
% respect to the bare mass~$m_0$. Some diagrams then produce the $1/\epsilon$
% pole as the result of the infrared (IR) divergence. We will find that, quite
% interestingly, these $1/\epsilon$ poles, in conjunction with the $1/\epsilon$
% poles arising from the UV divergence of other diagrams, are cancelled out
% in~$c_i(t)$; the $D\to4$ limit of~$c_i(t)$ is finite.

% \begin{figure}
% \centering
% \includegraphics[width=3cm,clip]{images/A01}
% \caption{A01}
% \end{figure}

% \begin{figure}
% \centering
% \includegraphics[width=3cm,clip]{images/A02}
% \caption{A02}
% \end{figure}

\begin{figure}
\begin{minipage}{0.3\textwidth}
\begin{center}
\includegraphics[width=3cm,clip]{images/A03}
\caption{A03}
\label{fig:14}
\end{center}
\end{minipage}
\begin{minipage}{0.3\textwidth}
\begin{center}
\includegraphics[width=3cm,clip]{images/A04}
\caption{A04}
\label{fig:15}
\end{center}
\end{minipage}
\begin{minipage}{0.3\textwidth}
\begin{center}
\includegraphics[width=3cm,clip]{images/A05}
\caption{A05}
\label{fig:16}
\end{center}
\end{minipage}
\end{figure}

\begin{figure}
\begin{minipage}{0.3\textwidth}
\begin{center}
\includegraphics[width=3cm,clip]{images/A06}
\caption{A06}
\label{fig:17}
\end{center}
\end{minipage}
\begin{minipage}{0.3\textwidth}
\begin{center}
\includegraphics[width=3cm,clip]{images/A07}
\caption{A07}
\label{fig:18}
\end{center}
\end{minipage}
\begin{minipage}{0.3\textwidth}
\begin{center}
\includegraphics[width=3cm,clip]{images/A08}
\caption{A08}
\label{fig:19}
\end{center}
\end{minipage}
\end{figure}

\begin{figure}
\begin{minipage}{0.3\textwidth}
\begin{center}
\includegraphics[width=3cm,clip]{images/A09}
\caption{A09}
\label{fig:20}
\end{center}
\end{minipage}
\begin{minipage}{0.3\textwidth}
\begin{center}
\includegraphics[width=3cm,clip]{images/A10}
\caption{A10}
\label{fig:21}
\end{center}
\end{minipage}
\begin{minipage}{0.3\textwidth}
\begin{center}
\includegraphics[width=3cm,clip]{images/A11}
\caption{A11}
\label{fig:22}
\end{center}
\end{minipage}
\end{figure}

\begin{figure}
\begin{minipage}{0.3\textwidth}
\begin{center}
\includegraphics[width=3cm,clip]{images/A12}
\caption{A12}
\label{fig:23}
\end{center}
\end{minipage}
\begin{minipage}{0.3\textwidth}
\begin{center}
\includegraphics[width=3cm,clip]{images/A13}
\caption{A13}
\label{fig:24}
\end{center}
\end{minipage}
\begin{minipage}{0.3\textwidth}
\begin{center}
\includegraphics[width=3cm,clip]{images/A14}
\caption{A14}
\label{fig:25}
\end{center}
\end{minipage}
\end{figure}

\begin{figure}
\begin{minipage}{0.3\textwidth}
\begin{center}
\includegraphics[width=3cm,clip]{images/A15}
\caption{A15}
\label{fig:26}
\end{center}
\end{minipage}
\begin{minipage}{0.3\textwidth}
\begin{center}
\includegraphics[width=3cm,clip]{images/A16}
\caption{A16}
\label{fig:27}
\end{center}
\end{minipage}
\begin{minipage}{0.3\textwidth}
\begin{center}
\includegraphics[width=3cm,clip]{images/A17}
\caption{A17}
\label{fig:28}
\end{center}
\end{minipage}
\end{figure}

\begin{figure}
\begin{minipage}{0.3\textwidth}
\begin{center}
\includegraphics[width=3cm,clip]{images/A18}
\caption{A18}
\label{fig:29}
\end{center}
\end{minipage}
\begin{minipage}{0.3\textwidth}
\begin{center}
\includegraphics[width=3cm,clip]{images/A19}
\caption{A19}
\label{fig:30}
\end{center}
\end{minipage}
\end{figure}
For $\zeta_{1j}^{(1)}(t)$, diagrams A03, A04, A05, A06, A07, A08, A09, A11, A13,
A14, A15, A16, and A19 in~Figs.~\ref{fig:14}--\ref{fig:30}
contribute.\footnote{Diagrams A10, A12, A17 (where the dotted line represents
the ghost propagator), and A18 in these figures correspond to the conventional
wave function renormalization and should be omitted in computing the operator
mixing.} In these and following diagrams, the cross generically represents one
of composite operators in~Eqs.~\eqref{eq:(4.1)}--\eqref{eq:(4.5)} [in the
present case, $\Tilde{\mathcal{O}}_{1\mu\nu}(t,x)$]. Apart from~A19, we can
borrow the results from~Ref.~\cite{Suzuki:2013gza} for these diagrams. For
completeness, these results are reproduced in the present convention
in~Table~\ref{table:0} of~Appendix~\ref{sec:C}. Combined with the contribution
of diagram~A19, we have
\begin{align}
   \zeta_{11}^{(1)}(t)
   &=\frac{g_0^2}{(4\pi)^2}C_2(G)
%   \left[\frac{11}{3}\epsilon(t)^{-1}+\frac{77}{24}\right],
   \left[\frac{11}{3}\epsilon(t)^{-1}+\frac{7}{3}\right],
\label{eq:(4.32)}\\
   \zeta_{12}^{(1)}(t)
   &=\frac{g_0^2}{(4\pi)^2}C_2(G)
%   \left[-\frac{11}{12}\epsilon(t)^{-1}-\frac{1}{8}\right],
   \left[-\frac{11}{12}\epsilon(t)^{-1}-\frac{1}{6}\right],
\label{eq:(4.33)}\\
   \zeta_{13}^{(1)}(t)
   &=\frac{g_0^4}{(4\pi)^2}C_2(R)
   \left[-\frac{2}{3}\epsilon(t)^{-1}-\frac{7}{18}\right],
\label{eq:(4.34)}\\
   \zeta_{14}^{(1)}(t)
   &=\frac{g_0^4}{(4\pi)^2}C_2(R)
   \left[\frac{1}{3}\epsilon(t)^{-1}-\frac{7}{18}\right],
\label{eq:(4.35)}\\
   \zeta_{15}^{(1)}(t)
   &=\frac{g_0^4}{(4\pi)^2}C_2(R)
   \left[3\epsilon(t)^{-1}+\frac{5}{2}\right],
\label{eq:(4.36)}
\end{align}
where
\begin{equation}
   \epsilon(t)^{-1}\equiv\frac{1}{\epsilon}+\ln(8\pi t).
\label{eq:(4.37)}
\end{equation}
By considering the trace part of~$\Tilde{\mathcal{O}}_{1\mu\nu}$, from these we
further have
\begin{align}
   \zeta_{21}^{(1)}(t)
   &=0,
\label{eq:(4.38)}\\
   \zeta_{22}^{(1)}(t)
   &=\zeta_{11}^{(1)}(t)+\zeta_{12}^{(1)}(t)(4-2\epsilon)
%   =\frac{g_0^2}{(4\pi)^2}C_2(G)\frac{109}{24},
   =\frac{g_0^2}{(4\pi)^2}C_2(G)\frac{7}{2},
\label{eq:(4.39)}\\
   \zeta_{23}^{(1)}(t)
   &=0,
\label{eq:(4.40)}\\
   \zeta_{24}^{(1)}(t)
   &=2\zeta_{13}^{(1)}(t)+\zeta_{14}^{(1)}(t)(4-2\epsilon)
   =\frac{g_0^4}{(4\pi)^2}C_2(R)(-3),
\label{eq:(4.41)}\\
   \zeta_{25}^{(1)}(t)
   &=\zeta_{15}^{(1)}(t)(4-2\epsilon)
   =\frac{g_0^4}{(4\pi)^2}C_2(R)
   \left[12\epsilon(t)^{-1}+4\right].
\label{eq:(4.42)}
\end{align}

% \begin{figure}
% \centering
% \includegraphics[width=3cm,clip]{images/B01}
% \caption{B01}
% \end{figure}

% \begin{figure}
% \centering
% \includegraphics[width=3cm,clip]{images/B02}
% \caption{B02}
% \end{figure}

\begin{figure}
\begin{minipage}{0.3\textwidth}
\begin{center}
\includegraphics[width=3cm,clip]{images/B03}
\caption{B03}
\label{fig:31}
\end{center}
\end{minipage}
\begin{minipage}{0.3\textwidth}
\begin{center}
\includegraphics[width=3cm,clip]{images/B04}
\caption{B04}
\label{fig:32}
\end{center}
\end{minipage}
\begin{minipage}{0.3\textwidth}
\begin{center}
\includegraphics[width=3cm,clip]{images/B05}
\caption{B05}
\label{fig:33}
\end{center}
\end{minipage}
\end{figure}

\begin{figure}
\begin{minipage}{0.3\textwidth}
\begin{center}
\includegraphics[width=3cm,clip]{images/B06}
\caption{B06}
\label{fig:34}
\end{center}
\end{minipage}
\begin{minipage}{0.3\textwidth}
\begin{center}
\includegraphics[width=3cm,clip]{images/B07}
\caption{B07}
\label{fig:35}
\end{center}
\end{minipage}
\begin{minipage}{0.3\textwidth}
\begin{center}
\includegraphics[width=3cm,clip]{images/B08}
\caption{B08}
\label{fig:36}
\end{center}
\end{minipage}
\end{figure}

\begin{figure}
\begin{minipage}{0.3\textwidth}
\begin{center}
\includegraphics[width=3cm,clip]{images/B09}
\caption{B09}
\label{fig:37}
\end{center}
\end{minipage}
\begin{minipage}{0.3\textwidth}
\begin{center}
\includegraphics[width=3cm,clip]{images/B10}
\caption{B10}
\label{fig:38}
\end{center}
\end{minipage}
\begin{minipage}{0.3\textwidth}
\begin{center}
\includegraphics[width=3cm,clip]{images/B11}
\caption{B11}
\label{fig:39}
\end{center}
\end{minipage}
\end{figure}

\begin{figure}
\begin{minipage}{0.3\textwidth}
\begin{center}
\includegraphics[width=3cm,clip]{images/B12}
\caption{B12}
\label{fig:40}
\end{center}
\end{minipage}
\begin{minipage}{0.3\textwidth}
\begin{center}
\includegraphics[width=3cm,clip]{images/B13}
\caption{B13}
\label{fig:41}
\end{center}
\end{minipage}
\begin{minipage}{0.3\textwidth}
\begin{center}
\includegraphics[width=3cm,clip]{images/B14}
\caption{B14}
\label{fig:42}
\end{center}
\end{minipage}
\end{figure}

\begin{figure}
\begin{minipage}{0.3\textwidth}
\begin{center}
\includegraphics[width=3cm,clip]{images/B15}
\caption{B15}
\label{fig:43}
\end{center}
\end{minipage}
\begin{minipage}{0.3\textwidth}
\begin{center}
\includegraphics[width=3cm,clip]{images/B16}
\caption{B16}
\label{fig:44}
\end{center}
\end{minipage}
\begin{minipage}{0.3\textwidth}
\begin{center}
\includegraphics[width=3cm,clip]{images/B17}
\caption{B17}
\label{fig:45}
\end{center}
\end{minipage}
\end{figure}

\begin{figure}
\begin{minipage}{0.3\textwidth}
\begin{center}
\includegraphics[width=3cm,clip]{images/B18}
\caption{B18}
\label{fig:46}
\end{center}
\end{minipage}
\end{figure}

For $\zeta_{3j}^{(1)}(t)$ with $j=1$ and~$2$, diagrams B03, B04, B08, B09, B10,
B11, and~B12 in~Figs.~\ref{fig:31}--\ref{fig:40} contribute. The contribution
of each diagram is tabulated in~Table~\ref{table:2}. For $\zeta_{3j}^{(1)}(t)$
with $j=3$, $4$, and~$5$, diagrams B06, B07, B13, B14, B15, B16, B17, and B18
in~Figs.~\ref{fig:34}--\ref{fig:46} contribute and their contributions are
tabulated in~Table~\ref{table:3}.
%
\begin{table}
\caption{$\zeta_{3j}^{(1)}$ in units of $\frac{1}{(4\pi)^2}T(R)N_{\mathrm{f}}$}
\label{table:2}
\begin{center}
\renewcommand{\arraystretch}{2.2}
\setlength{\tabcolsep}{20pt}
\begin{tabular}{crr}
\toprule
 diagram & \multicolumn{1}{c}{$\zeta_{31}^{(1)}(t)$}
 & \multicolumn{1}{c}{$\zeta_{32}^{(1)}(t)$} \\
\midrule
B03  & $-\dfrac{16}{3}\epsilon(t)^{-1}-\dfrac{64}{9}$ & $\dfrac{4}{3}\epsilon(t)^{-1}+\dfrac{25}{9}$ \\
B04  & $0$ & $0$ \\
B08  & $0$ & $-\dfrac{1}{3}$ \\
B09  & $-\dfrac{5}{12}$ & $-\dfrac{23}{144}$ \\
B10  & $\dfrac{5}{12}$ & $-\dfrac{1}{16}$ \\
B11  & $\dfrac{10}{9}$ & $\dfrac{7}{9}$ \\
B12  & $0$ & $0$ \\
\bottomrule
\end{tabular}
\end{center}
\end{table}
%
\begin{table}
\caption{$\zeta_{3j}^{(1)}$ in units of $\frac{g_0^2}{(4\pi)^2}C_2(R)$}
\label{table:3}
\begin{center}
\renewcommand{\arraystretch}{2.2}
\setlength{\tabcolsep}{20pt}
\begin{tabular}{crrr}
\toprule
 diagram & \multicolumn{1}{c}{$\zeta_{33}^{(1)}(t)$}
 & \multicolumn{1}{c}{$\zeta_{34}^{(1)}(t)$}
 & \multicolumn{1}{c}{$\zeta_{35}^{(1)}(t)$} \\
\midrule
 B06 & $-\dfrac{1}{3}\epsilon(t)^{-1}+\dfrac{1}{18}$ & $\dfrac{2}{3}\epsilon(t)^{-1}+\dfrac{5}{9}$ & $8\epsilon(t)^{-1}+8$ \\
 B07 & $2\epsilon(t)^{-1}+2$ & $-2\epsilon(t)^{-1}-2$ & $-8\epsilon(t)^{-1}-8$ \\
 B13 & $0$ & $1$ & $0$ \\
 B14 & $0$ & $0$ & $0$ \\
 B15 & $2\epsilon(t)^{-1}+2$ & $0$ & $0$ \\
 B16 & $0$ & $0$ & $0$ \\
 B17 & $-2$ & $0$ & $0$ \\
 B18 & $-4\epsilon(t)^{-1}-2$ & $0$ & $0$ \\
\bottomrule
\end{tabular}
\end{center}
\end{table}
%
As the sum of these contributions, we have
\begin{align}
   \zeta_{31}^{(1)}(t)
   &=\frac{1}{(4\pi)^2}T(R)N_{\mathrm{f}}
   \left[-\frac{16}{3}\epsilon(t)^{-1}-6\right],
\label{eq:(4.43)}\\
   \zeta_{32}^{(1)}(t)
   &=\frac{1}{(4\pi)^2}T(R)N_{\mathrm{f}}
   \left[\frac{4}{3}\epsilon(t)^{-1}+3\right],
\label{eq:(4.44)}\\
   \zeta_{33}^{(1)}(t)
   &=\frac{g^2}{(4\pi)^2}C_2(R)
   \left[3\frac{1}{\epsilon}-\Phi(t)\right]
   +\frac{g_0^2}{(4\pi)^2}C_2(R)
   \left[-\frac{1}{3}\epsilon(t)^{-1}+\frac{1}{18}\right],
\label{eq:(4.45)}\\
   \zeta_{34}^{(1)}(t)
   &=\frac{g_0^2}{(4\pi)^2}C_2(R)
   \left[-\frac{4}{3}\epsilon(t)^{-1}-\frac{4}{9}\right],
\label{eq:(4.46)}\\
   \zeta_{35}^{(1)}(t)
   &=0,
\label{eq:(4.47)}
\end{align}
where in $\zeta_{33}^{(1)}(t)$~\eqref{eq:(4.45)} the first term in the
right-hand side comes from the conversion from the un-ringed fields to the
ringed fields in~Eqs.~\eqref{eq:(3.20)} and~\eqref{eq:(3.21)}---recall
Eq.~\eqref{eq:(3.25)}; the combination~$\Phi(t)$ is given
by~Eq.~\eqref{eq:(3.27)}. From these, we further have
\begin{align}
   \zeta_{41}^{(1)}(t)
   &=0,
\label{eq:(4.48)}\\
   \zeta_{42}^{(1)}(t)
   &=\frac{1}{2}\zeta_{31}^{(1)}(t)+\frac{1}{2}\zeta_{32}^{(1)}(t)(4-2\epsilon)
   =\frac{1}{(4\pi)^2}T(R)N_{\mathrm{f}}
   \frac{5}{3},
\label{eq:(4.49)}\\
   \zeta_{43}^{(1)}(t)
   &=0,
\label{eq:(4.50)}\\
   \zeta_{44}^{(1)}(t)
   &=\zeta_{33}^{(1)}(t)+\frac{1}{2}\zeta_{34}^{(1)}(t)(4-2\epsilon)
\notag\\
   &=\frac{g^2}{(4\pi)^2}C_2(R)
   \left[3\frac{1}{\epsilon}-\Phi(t)\right]
   +\frac{g_0^2}{(4\pi)^2}C_2(R)
   \left[-3\epsilon(t)^{-1}+\frac{1}{2}\right],
\label{eq:(4.51)}\\
   \zeta_{45}^{(1)}(t)
   &=0.
\label{eq:(4.52)}
\end{align}

Finally,
\begin{equation}
   \zeta_{51}^{(1)}(t)
   =\zeta_{52}^{(1)}(t)=\zeta_{53}^{(1)}(t)=\zeta_{54}^{(1)}(t)=0,
\label{eq:(4.53)}
\end{equation}
and $\zeta_{55}^{(1)}(t)$ is given by the sum of the contributions of one-loop
diagrams in~Table~\ref{table:4} and the conversion factor to the ringed fields:
%
\begin{table}
\caption{$\zeta_{55}^{(1)}(t)$ in units of~$\frac{g_0^2}{(4\pi)^2}C_2(R)$.}
\label{table:4}
\begin{center}
\renewcommand{\arraystretch}{2.2}
\setlength{\tabcolsep}{20pt}
\begin{tabular}{cr}
\toprule
 diagram & \multicolumn{1}{c}{$\zeta_{55}^{(1)}(t)$} \\
\midrule
 B06 & $-4\epsilon(t)^{-1}-2$ \\
 B13 & $0$ \\
 B14 & $0$ \\
 B15 & $2\epsilon(t)^{-1}+2$ \\
 B16 & $0$ \\
 B18 & $-4\epsilon(t)^{-1}-2$ \\
\bottomrule
\end{tabular}
\end{center}
\end{table}
%
\begin{align}
   \zeta_{55}^{(1)}(t)
   =\frac{g^2}{(4\pi)^2}C_2(R)
   \left[6\frac{1}{\epsilon}-\Phi(t)\right]
   +\frac{g_0^2}{(4\pi)^2}C_2(R)
   \left[-6\epsilon(t)^{-1}-2\right].
\label{eq:(4.54)}
\end{align}

We have now obtained all $\zeta_{ij}^{(1)}(t)$ in~Eqs.~\eqref{eq:(4.24)}
and~\eqref{eq:(4.25)}. Then, since the matrix~$\zeta_{ij}(t)$ in the tree-level
approximation is a unit matrix, it is straightforward to invert the
matrix~$\zeta_{ij}(t)$ in the one-loop approximation; Eqs.~\eqref{eq:(4.15)}
and~\eqref{eq:(4.16)} thus yield
\begin{align}
   c_1(t)&=\frac{1}{g_0^2}
   \left[1-\zeta_{11}^{(1)}(t)\right]
   -\frac{1}{4}\zeta_{31}^{(1)}(t),
\label{eq:(4.55)}\\
   c_2(t)&=\frac{1}{g_0^2}
   \left(-\frac{1}{2}\epsilon\right)\zeta_{12}^{(1)}(t)
   +\left(\frac{3}{4}-\frac{1}{2}\epsilon\right)\zeta_{32}^{(1)}(t)
   +\frac{3}{16}\zeta_{31}^{(1)}(t),
\label{eq:(4.56)}\\
   c_3(t)&=\left\{\frac{1}{g_0^2}
   \left[-\zeta_{13}^{(1)}(t)\right]
   +\frac{1}{4}-\frac{1}{4}\zeta_{33}^{(1)}(t)\right\}Z(\epsilon)^{-1},
\label{eq:(4.57)}\\
   c_4(t)&=\left\{\frac{1}{g_0^2}
   \left[-\frac{1}{2}\epsilon\zeta_{14}^{(1)}(t)
   -\frac{3}{2}\zeta_{13}^{(1)}(t)\right]
   +\left(\frac{3}{4}-\frac{1}{2}\epsilon\right)\zeta_{34}^{(1)}(t)\right\}
   Z(\epsilon)^{-1},
\label{eq:(4.58)}\\
   c_5(t)&=\left\{\frac{1}{g_0^2}
   \left[-\frac{1}{2}\epsilon\zeta_{15}^{(1)}(t)\right]
   -1+\zeta_{55}^{(1)}(t)\right\}Z(\epsilon)^{-1}.
\label{eq:(4.59)}
\end{align}
Then, by using the renormalized gauge coupling in the $\text{MS}$
scheme~\eqref{eq:(A1)}, to the one-loop order, we have (for~$\epsilon\to0$)
\begin{align}
   c_1(t)&
%   =\frac{1}{g^2}-b_0\ln(8\pi\mu^2t)-\frac{7}{8}\frac{1}{(4\pi)^2}
%   \left[\frac{11}{3}C_2(G)
%   -\frac{12}{7}T(R)N_{\mathrm{f}}\right],
   =\frac{1}{g^2}-b_0\ln(8\pi\mu^2t)-\frac{1}{(4\pi)^2}
   \left[\frac{7}{3}C_2(G)
   -\frac{3}{2}T(R)N_{\mathrm{f}}\right],
\label{eq:(4.60)}\\
   c_2(t)&
   =\frac{1}{8}
   \frac{1}{(4\pi)^2}
   \left[\frac{11}{3}C_2(G)
   +\frac{11}{3}T(R)N_{\mathrm{f}}\right],
\label{eq:(4.61)}\\
   c_3(t)&
   =\frac{1}{4}\left\{1+\frac{g^2}{(4\pi)^2}C_2(R)
   \left[\frac{3}{2}+\ln(432)\right]\right\},
\label{eq:(4.62)}\\
   c_4(t)&=\frac{1}{8}d_0g^2,
\label{eq:(4.63)}\\
   c_5(t)&=-\left\{1+\frac{g^2}{(4\pi)^2}C_2(R)
   \left[3\ln(8\pi\mu^2t)+\frac{7}{2}+\ln(432)\right]\right\}.
\label{eq:(4.64)}
\end{align}
Since the $c_i(t)$ in~Eq.~\eqref{eq:(4.14)} connect the finite energy--momentum
tensor and UV-finite local products $\Tilde{\mathcal{O}}_{i\mu\nu}(t,x)$
constructed from (ringed) flowed fields, they should be UV finite. That our
explicit one-loop calculation of~$c_i(t)$ confirms this finiteness is quite
reassuring.

If one prefers the $\overline{\text{MS}}$ scheme instead of the $\text{MS}$
scheme assumed in above expressions, it suffices to make the replacement
\begin{equation}
    \mu^2\to\frac{\mathrm{e}^{\gamma_{\mathrm{E}}}}{4\pi}\mu^2,
\label{eq:(4.65)}
\end{equation}
where $\gamma_{\mathrm{E}}$ is Euler's constant.

\subsection{A consistency check: The trace anomaly}
It is interesting to see that Eq.~\eqref{eq:(4.14)} with $c_i(t)$
in~Eqs.~\eqref{eq:(4.60)}--\eqref{eq:(4.64)} in fact reproduces the trace
anomaly~\eqref{eq:(2.11)} in the one-loop approximation. At first brief glance,
$c_2(t)$ in~Eq.~\eqref{eq:(4.61)} is incompatible with the correct trace
anomaly, because $\delta_{\mu\nu}[\Tilde{\mathcal{O}}_{1\mu\nu}
-(1/4)\Tilde{\mathcal{O}}_{2\mu\nu}]=0$ for~$D=4$ and
$\delta_{\mu\nu}\Tilde{\mathcal{O}}_{2\mu\nu}
=4\left\{F_{\rho\sigma}^aF_{\rho\sigma}^a\right\}_R(x)+O(g^2)$ [the last equality
follows from Eq.~\eqref{eq:(A17)}]. On the other hand, $4c_2(t)$
from~Eq.~\eqref{eq:(4.61)} is not identical to~$b_0/2$, where $b_0$ is the
first coefficient of the beta function~\eqref{eq:(2.16)}, the correct one-loop
coefficient of the trace anomaly.

This is a premature judgment, however. In fact, $\zeta_{42}^{(1)}(t)$
in~Eq.~\eqref{eq:(4.49)} shows that there exists an operator mixing of the form
\begin{equation}
   \mathring{\Bar{\chi}}(t,x)
   \overleftrightarrow{\Slash{D}}
   \mathring{\chi}(t,x)
   =\left\{\Bar{\psi}\overleftrightarrow{\Slash{D}}\psi\right\}_R(x)
   +\frac{1}{(4\pi)^2}T(R)N_{\mathrm{f}}\frac{5}{3}
   \left\{F_{\rho\sigma}^aF_{\rho\sigma}^a\right\}_R(x)
   +\dotsb.
\label{eq:(4.66)}
\end{equation}
Then the last term precisely fills the difference between $4c_2(t)$
and~$b_0/2$.

In this way, to the one-loop order, we have
\begin{align}
   \delta_{\mu\nu}\left\{T_{\mu\nu}\right\}_R(x)
   &=\frac{1}{2}\frac{1}{(4\pi)^2}
   \left[\frac{11}{3}C_2(G)-\frac{4}{3}T(R)N_{\mathrm{f}}\right]
   \left\{F_{\rho\sigma}^aF_{\rho\sigma}^a\right\}_R(x)
\notag\\
   &\qquad{}
   -\frac{3}{2}
   \left\{\Bar{\psi}\overleftrightarrow{\Slash{D}}\psi\right\}_R(x)
   -\left[4+\frac{g^2}{(4\pi)^2}6C_2(R)\right]
   m\left\{\Bar{\psi}\psi\right\}_R(x).
\label{eq:(4.67)}
\end{align}
This reproduces the trace anomaly~\eqref{eq:(2.11)} in the one-loop level if
one uses the equation of motion of renormalized field,
\begin{equation}
   \left\{\Bar{\psi}\overleftrightarrow{\Slash{D}}\psi\right\}_R(x)
   =-2m\left\{\Bar{\psi}\psi\right\}_R(x),
\label{eq:(4.68)}
\end{equation}
whose use is justified when there is no other operator in the point~$x$ as we
are assuming. We observe that our one-loop result
in~Eqs.~\eqref{eq:(4.60)}--\eqref{eq:(4.64)} is consistent with the trace
anomaly.

In Ref.~\cite{Suzuki:2013gza}, for the pure Yang--Mills theory, the
next-to-leading (two-loop order) term in~$c_2(t)$ was determined as
\begin{equation}
   c_2(t)=\frac{1}{8}b_0
   -\frac{1}{8}b_0
%   \left[\frac{1}{(4\pi)^2}G_2(G)\frac{109}{24}-\frac{b_1}{b_0}\right]g^2,
   \left[\frac{1}{(4\pi)^2}G_2(G)\frac{7}{2}-\frac{b_1}{b_0}\right]g^2,
\label{eq:(4.69)}
\end{equation}
where $b_0=[1/(4\pi)^2](11/3)C_2(G)$ and~$b_1=[1/(4\pi)^4](34/3)C_2(G)^2$, by
imposing that the expression~\eqref{eq:(4.14)} reproduces the trace
anomaly~\eqref{eq:(2.11)} to the two-loop order. For the present system with
fermions, however, it seems that this requirement alone cannot fix the
next-to-leading terms in~$c_2(t)$ and in~$c_4(t)$; so we are content with the
one-loop formulas, Eqs.~\eqref{eq:(4.60)}--\eqref{eq:(4.64)}, in the present
paper treating a system containing fermions.

\subsection{Master formula}
From~Eq.~\eqref{eq:(4.14)}, the energy--momentum tensor is given by the $t\to0$
limit,
\begin{align}
   \left\{T_{\mu\nu}\right\}_R(x)
   &=\lim_{t\to0}\biggl\{c_1(t)\left[
   \Tilde{\mathcal{O}}_{1\mu\nu}(t,x)
   -\frac{1}{4}\Tilde{\mathcal{O}}_{2\mu\nu}(t,x)
   \right]
\notag\\
   &\qquad\qquad{}
   +c_2(t)\left[
   \Tilde{\mathcal{O}}_{2\mu\nu}(t,x)
   -\left\langle\Tilde{\mathcal{O}}_{2\mu\nu}(t,x)\right\rangle
   \right]
\notag\\
   &\qquad\qquad\qquad{}
   +c_3(t)\left[
   \Tilde{\mathcal{O}}_{3\mu\nu}(t,x)
   -2\Tilde{\mathcal{O}}_{4\mu\nu}(t,x)
   -\left\langle
   \Tilde{\mathcal{O}}_{3\mu\nu}(t,x)
   -2\Tilde{\mathcal{O}}_{4\mu\nu}(t,x)
   \right\rangle
   \right]
\notag\\
   &\qquad\qquad\qquad\qquad{}
   +c_4(t)\left[
   \Tilde{\mathcal{O}}_{4\mu\nu}(t,x)
   -\left\langle\Tilde{\mathcal{O}}_{4\mu\nu}(t,x)\right\rangle
   \right]
\notag\\
   &\qquad\qquad\qquad\qquad\qquad{}
   +c_5(t)\left[
   \Tilde{\mathcal{O}}_{5\mu\nu}(t,x)
   -\left\langle\Tilde{\mathcal{O}}_{5\mu\nu}(t,x)\right\rangle
   \right]\biggr\},
\label{eq:(4.70)}
\end{align}
where operators in the right-hand side are given
by~Eqs.~\eqref{eq:(4.1)}--\eqref{eq:(4.5)}. One may further use the identities
\begin{equation}
   2\left\langle\Tilde{\mathcal{O}}_{3\mu\nu}(t,x)\right\rangle
   =\left\langle\Tilde{\mathcal{O}}_{4\mu\nu}(t,x)\right\rangle
   =\frac{-2\dim(R)N_{\mathrm{f}}}{(4\pi)^2t^2}\delta_{\mu\nu}
\label{eq:(4.71)}
\end{equation}
to make the expression a little simpler. Applying the consequence of the
renormalization group argument, Eqs.~\eqref{eq:(4.22)}--\eqref{eq:(4.23)},
to~Eqs.~\eqref{eq:(4.60)}--\eqref{eq:(4.64)}, we have
\begin{align}
   c_1(t)&
%   =\frac{1}{\Bar{g}(1/\sqrt{8t})^2}
%   -b_0\ln\pi-\frac{7}{8}\frac{1}{(4\pi)^2}
%   \left[\frac{11}{3}C_2(G)
%   -\frac{12}{7}T(R)N_{\mathrm{f}}\right],
   =\frac{1}{\Bar{g}(1/\sqrt{8t})^2}
   -b_0\ln\pi-\frac{1}{(4\pi)^2}
   \left[\frac{7}{3}C_2(G)
   -\frac{3}{2}T(R)N_{\mathrm{f}}\right],
\label{eq:(4.72)}\\
   c_2(t)&
   =\frac{1}{8}
   \frac{1}{(4\pi)^2}
   \left[\frac{11}{3}C_2(G)
   +\frac{11}{3}T(R)N_{\mathrm{f}}\right],
\label{eq:(4.73)}\\
   c_3(t)&
   =\frac{1}{4}\left\{1+\frac{\Bar{g}(1/\sqrt{8t})^2}{(4\pi)^2}C_2(R)
   \left[\frac{3}{2}+\ln(432)\right]\right\},
\label{eq:(4.74)}\\
   c_4(t)&=\frac{1}{8}d_0\Bar{g}(1/\sqrt{8t})^2,
\label{eq:(4.75)}\\
   c_5(t)&=-\frac{\Bar{m}(1/\sqrt{8t})}{m}
   \left\{1+\frac{\Bar{g}(1/\sqrt{8t})^2}{(4\pi)^2}C_2(R)
   \left[3\ln\pi+\frac{7}{2}+\ln(432)\right]\right\},
\label{eq:(4.76)}
\end{align}
where $\Bar{g}(q)$ is the running gauge coupling in the $\text{MS}$ scheme.
For going from the $\text{MS}$ scheme to the $\overline{\text{MS}}$ scheme, it
suffices to make the following replacement corresponding
to~Eq.~\eqref{eq:(4.65)}, 
\begin{equation}
   \ln\pi\to\gamma_{\mathrm{E}}-2\ln2,
\label{eq:(4.77)}
\end{equation}
in the above expressions. Equation~\eqref{eq:(4.70)} with
Eqs.~\eqref{eq:(4.72)}--\eqref{eq:(4.76)} is our main result. Note that the
renormalized mass parameter~$m$ in~$c_5(t)$~\eqref{eq:(4.76)} and that
in~$\Tilde{\mathcal{O}}_{5\mu\nu}(t,x)$~\eqref{eq:(4.5)} are redundant
in~Eq.~\eqref{eq:(4.70)} because they are cancelled out in the product. For the
running mass parameter~$\Bar{m}(1/\sqrt{8t})$ in~$c_5(t)$ for~$t\to0$, one may
use the relation
\begin{align}
   \Bar{m}(q)&=[2b_0\Bar{g}(q)^2]^{d_0/2b_0}
   \exp\left\{\int_0^{\Bar{g}(q)}\mathrm{d}g\,
   \left[-\frac{\gamma_m(g)}{\beta(g)}-\frac{d_0}{b_0g}\right]\right\}m^\infty
\notag\\
   &=\left(\frac{2}{\ell}\right)^{d_0/2b_0}
   \left[1-\frac{d_0b_1}{2b_0^3\ell}(1+\ln\ell)+\frac{d_1}{2b_0^2\ell}
   +O(\ell^{-2})\right]m^\infty,\qquad\ell\equiv\ln(q^2/\Lambda^2),
\label{eq:(4.78)}
\end{align}
where $m^\infty$ denotes the renormalization group invariant mass. For the
massive fermion, $m^\infty$~may be determined by using the method established
in~Ref.~\cite{Capitani:1998mq}, for example. For the massless fermion,
$m^\infty=0$ and we can simply discard the last line of~Eq.~\eqref{eq:(4.70)}.
